
\documentclass[11pt]{article}
\usepackage[margin=1in]{geometry}
\usepackage{amsmath,amssymb}
\usepackage{enumitem}
\usepackage[T1]{fontenc}
\usepackage{lmodern}
\setlength{\parskip}{0.5em}
\setlength{\parindent}{0pt}

\begin{document}

\begin{center}
{\LARGE \textbf{IB Physics 2016--2020}}\\
{\large \textbf{Mapped Collection for D2 + D3}}\\
\vspace{0.3em}
\textit{Faithful paraphrases of old-syllabus questions matched to the modern D2/D3 field units.}
\end{center}

\textbf{Mapping used.} Since 2016--2020 papers predate the 2025 syllabus, this collection maps modern D2/D3 to the closest old-syllabus material on electric fields, electric potential, magnetic forces, motion of charges in fields, and related electromagnetic field models.

\section*{Problems}

\begin{enumerate}[leftmargin=*, itemsep=1.2em]

\item \textbf{M16 HL Paper 2, Q6(b) --- conductor moving charges in a magnetic field}\\
A metallic conductor of width $L$ contains charge carriers moving with speed $v$. The conductor is placed in a uniform magnetic field $B$ directed into the page.
\begin{enumerate}[label=(\alph*)]
\item Show that the potential difference established across the conductor is $V=vBL$.
\item State the physical origin of this potential difference in terms of magnetic force on the carriers and charge separation.
\end{enumerate}

\item \textbf{M17 TZ2 HL Paper 2, Q6(a) --- current in a conductor versus an insulator}\\
A power cable consists of many copper strands surrounded by insulating material. Both the copper and the insulator are placed in an electric field.
\begin{enumerate}[label=(\alph*)]
\item Explain, with reference to charge carriers, why a significant current exists in the copper but not in the insulator.
\item State the role of free electrons in the conductor and compare this with the charge behaviour in the insulating layer.
\end{enumerate}

\item \textbf{M17 TZ2 HL Paper 2, Q6(d) --- magnetic forces between parallel ac cables}\\
Two identical long cables are suspended parallel to one another and carry alternating current as part of a transmission system.
\begin{enumerate}[label=(\alph*)]
\item Explain qualitatively how the magnetic force between the cables changes during one full ac cycle.
\item State when the force is zero, when it is largest in magnitude, and whether it is attractive or repulsive.
\end{enumerate}

\item \textbf{M18 TZ2 HL Paper 2, Q8 --- thundercloud--Earth as a parallel-plate capacitor}\\
A negatively charged thundercloud base is modelled as one plate of a parallel-plate capacitor and the Earth as the other. The given data are:
\[
A=1.2\times 10^{8}\ \text{m}^2,\qquad Q=-25\ \text{C},\qquad d=1600\ \text{m},\qquad \varepsilon=8.8\times 10^{-12}\ \text{F m}^{-1}.
\]
\begin{enumerate}[label=(\alph*)]
\item Show that the capacitance is approximately $6.6\times 10^{-7}\,\text{F}$.
\item Calculate the potential difference between the thundercloud and the Earth.
\item Calculate the energy stored in the system.
\item The cloud discharges in a lightning strike. The time constant is $32\,\text{ms}$ and the strike lasts $18\,\text{ms}$. Show that roughly $11\,\text{C}$ of charge is transferred.
\item Hence estimate the average current during the strike.
\item State one assumption required in order to model the Earth--thundercloud system as a parallel-plate capacitor.
\end{enumerate}

\item \textbf{M19 TZ2 HL Paper 2, Q5 --- proton in a uniform magnetic field}\\
A proton moves on a circular path in a uniform magnetic field directed into the plane of the page.
\begin{enumerate}[label=(\alph*)]
\item On a diagram, indicate the direction of the magnetic force and the direction of the proton's velocity.
\item Given $v=2.16\times 10^6\,\text{m s}^{-1}$ and $B=0.042\,\text{T}$, calculate the radius of the circular path.
\item State why the path is circular and why the speed remains constant.
\end{enumerate}

\item \textbf{M19 TZ2 HL Paper 2, Q9(b) --- zero electric field point between two charges}\\
Two fixed charges $X$ and $Y$ have magnitudes $Q$ and $q$ respectively. A point $P$ lies on the line through the charges. The distance between $X$ and $Y$ is $0.600\,\text{m}$ and the distance between $P$ and $Y$ is $0.820\,\text{m}$. The electric field at $P$ is zero.
\begin{enumerate}[label=(\alph*)]
\item Determine the ratio $Q/q$ to one significant figure.
\item State what the field-line diagram implies about the signs of the charges.
\end{enumerate}

\item \textbf{N19 HL Paper 2, Q9 --- capacitor energy and redistribution of charge}\\
A system of capacitors is charged and then reconfigured.
\begin{enumerate}[label=(\alph*)]
\item Using capacitance and potential difference, calculate the initial energy stored in a capacitor.
\item Two capacitors with capacitances $18\,\mu\text{F}$ and $12\,\mu\text{F}$ share a total charge of $48\,\mu\text{C}$. Determine the charge on each capacitor.
\item Hence find the total final energy stored.
\item Explain why energy is lost when the capacitors are connected and charge redistributes.
\end{enumerate}

\item \textbf{N20 HL Paper 2, Q8 --- electric potential of a charged sphere}\\
A positively charged sphere produces a radial electric field.
\begin{enumerate}[label=(\alph*)]
\item Explain why the electric potential decreases as one moves radially outward from the sphere.
\item Sketch the variation of electric potential $V$ with distance $r$ from the centre, both inside and outside the sphere.
\item A particle of known charge has electric potential energy $1.7\times 10^{-16}\,\text{J}$ at a point where the electric potential is $1.1\times 10^3\,\text{V}$. Determine the charge on the particle.
\item Using the potential of a point charge, determine the charge on the sphere from suitable distance data.
\item State one reason physicists compare gravitational and electric field models.
\end{enumerate}

\end{enumerate}

\end{document}
